% GENERATED FILE. Edit canonical lessons, then run npm run generate:books.
% canonical-source-hash: fnv1a64:7da21dff4c0b7348
% canonical-lessons: FR-C172-plastique, FR-C172-coton, FR-C172-laine, FR-C172-soie, FR-C172-cuir

\chapter{Plastic, Cotton, Wool, Silk, Leather}
\label{ch:fr-plastic-cotton-wool-silk-leather}
\hlchaptermodality{172}

\begin{chapteropening}
\textbf{By the end of this chapter:} \emph{I can say plastic, cotton, wool, silk and leather.}

Complete the last lesson of chapter 172: I can say plastic, cotton, wool, silk and leather.
\end{chapteropening}

\section[le plastique]{le plastique — plastic}

\label{lesson:FR-C172-plastique}



\begin{quote}
Before the new one: say the French for a rock, then the French for a pebble.
\end{quote}

\subsection*{You'll want to know: le plastique}
\textbf{le plastique} — \textquotedblleft{}plastic\textquotedblright{}.

\begin{grammarlens}[title={What you've built}]
One more word for this chapter.
\end{grammarlens}

\subsection*{Your turn}
\begin{itemize}
\raggedright
  \item \emph{Say it:} \emph{le plastique}
  \item \emph{Say it:} \emph{le plastique}, once more
  \item \emph{Say it:} say \emph{le plastique}
  \item \emph{Recall:} say the French for a rock, then the French for a pebble, then say \emph{le plastique} again
\end{itemize}

\begin{tcolorbox}[breakable,colback=teal!4,colframe=teal!35!black,title={Before you move on}]
What does le plastique mean? (\textquotedblleft{}Plastic\textquotedblright{}.) Say it once more.
\end{tcolorbox}
\section[le coton]{le coton — cotton}

\label{lesson:FR-C172-coton}



\begin{quote}
Before the new one: say the French for a pebble, then the French for plastic.
\end{quote}

\subsection*{You'll want to know: le coton}
\textbf{le coton} — \textquotedblleft{}cotton\textquotedblright{}.

\begin{grammarlens}[title={What you've built}]
One more word for this chapter.
\end{grammarlens}

\subsection*{Your turn}
\begin{itemize}
\raggedright
  \item \emph{Say it:} \emph{le coton}
  \item \emph{Say it:} \emph{le coton}, once more
  \item \emph{Say it:} say \emph{le coton}
  \item \emph{Recall:} say the French for a pebble, then the French for plastic, then say \emph{le coton} again
\end{itemize}

\begin{tcolorbox}[breakable,colback=teal!4,colframe=teal!35!black,title={Before you move on}]
What does le coton mean? (\textquotedblleft{}Cotton\textquotedblright{}.) Say it once more.
\end{tcolorbox}
\section[la laine]{la laine — wool}

\label{lesson:FR-C172-laine}



\begin{quote}
Before the new one: say the French for plastic, then the French for cotton.
\end{quote}

\subsection*{You'll want to know: la laine}
\textbf{la laine} — \textquotedblleft{}wool\textquotedblright{}.

\begin{grammarlens}[title={What you've built}]
One more word for this chapter.
\end{grammarlens}

\subsection*{Your turn}
\begin{itemize}
\raggedright
  \item \emph{Say it:} \emph{la laine}
  \item \emph{Say it:} \emph{la laine}, once more
  \item \emph{Say it:} say \emph{la laine}
  \item \emph{Recall:} say the French for plastic, then the French for cotton, then say \emph{la laine} again
\end{itemize}

\begin{tcolorbox}[breakable,colback=teal!4,colframe=teal!35!black,title={Before you move on}]
What does la laine mean? (\textquotedblleft{}Wool\textquotedblright{}.) Say it once more.
\end{tcolorbox}
\section[la soie]{la soie — silk}

\label{lesson:FR-C172-soie}



\begin{quote}
Before the new one: say the French for cotton, then the French for wool.
\end{quote}

\subsection*{You'll want to know: la soie}
\textbf{la soie} — \textquotedblleft{}silk\textquotedblright{}.

\begin{grammarlens}[title={What you've built}]
One more word for this chapter.
\end{grammarlens}

\subsection*{Your turn}
\begin{itemize}
\raggedright
  \item \emph{Say it:} \emph{la soie}
  \item \emph{Say it:} \emph{la soie}, once more
  \item \emph{Say it:} say \emph{la soie}
  \item \emph{Recall:} say the French for cotton, then the French for wool, then say \emph{la soie} again
\end{itemize}

\begin{tcolorbox}[breakable,colback=teal!4,colframe=teal!35!black,title={Before you move on}]
What does la soie mean? (\textquotedblleft{}Silk\textquotedblright{}.) Say it once more.
\end{tcolorbox}
\section[le cuir]{le cuir — leather}

\label{lesson:FR-C172-cuir}



\begin{quote}
Before the new one: say the French for wool, then the French for silk.
\end{quote}

\subsection*{You'll want to know: le cuir}
\textbf{le cuir} — \textquotedblleft{}leather\textquotedblright{}.

\begin{grammarlens}[title={What you've built}]
One more word for this chapter.
\end{grammarlens}

\subsection*{Your turn}
\begin{itemize}
\raggedright
  \item \emph{Say it:} \emph{le cuir}
  \item \emph{Say it:} \emph{le cuir}, once more
  \item \emph{Say it:} say \emph{le cuir}
  \item \emph{Recall:} say the French for wool, then the French for silk, then say \emph{le cuir} again
\end{itemize}

\begin{tcolorbox}[breakable,colback=teal!4,colframe=teal!35!black,title={Before you move on}]
What does le cuir mean? (\textquotedblleft{}Leather\textquotedblright{}.) Say it once more.
\end{tcolorbox}
